\documentclass[12pt,a4paper]{article}
\usepackage[utf8]{inputenc}
\usepackage[french]{babel}
\usepackage[T1]{fontenc}
\usepackage{geometry}
\geometry{top=2.5cm, bottom=2.5cm, left=2.5cm, right=2.5cm}
\usepackage{amsmath, amssymb, amsfonts}
\usepackage{booktabs}
\usepackage{xcolor}
\usepackage{hyperref}
\usepackage{fancyhdr}

\hypersetup{
    colorlinks=true,
    linkcolor=blue,
    urlcolor=cyan,
    pdftitle={Documentation Decison Support - DAXX EU},
}

\pagestyle{fancy}
\fancyhf{}
\rhead{Aide à la Décision -- DAXX EU}
\lhead{Algorithmes, Notations, Méthodes, Risques et Précision}
\rfoot{Page \thepage}

\title{\textbf{Guide Complet du Module Decision Support \\ Algorithmes, Notations, Risques, Précision et Démonstration \\ Plateforme Décisionnelle -- DAXX EU}}
\author{\textbf{Département Modélisation \& Data Science -- DAXX EU}}
\date{\today}

\begin{document}

\maketitle
\tableofcontents
\newpage

\section{Sur quelles données se basent les calculs ?}
Le moteur décisionnel s'appuie sur une base de données de prix physiques réels extraite quotidiennement depuis la plateforme d'intelligence de marché **DAXX Pricing Intel**.

\begin{itemize}
    \item \textbf{Historique Temporel} : Les séries contiennent des relevés de prix quotidiens du 1er janvier 2021 à aujourd'hui (soit plus de 1300 jours d'historique).
    \item \textbf{Nature des Prix} : Il s'agit de prix spot intérieurs hors taxes (\textit{Domestic Spot Prices}) ou prix départ usine (\textit{Ex-Works}) exprimés en Yuans par tonne ($\text{\yen}/\text{tonne}$).
    \item \textbf{Alignement et Nettoyage} : Les jours fériés et week-ends chinois (périodes sans cotation) sont harmonisés par un alignement temporel strict. Les trous de données sont résolus par la méthode du \textbf{Forward-Fill} (le prix d'un jour férié est égal au dernier prix connu de la veille).
\end{itemize}

---

\section{Les 3 Méthodes d'Aide à la Décision (selon la disponibilité des données)}
Le moteur décisionnel applique automatiquement l'une des trois méthodes suivantes selon l'état et la densité des données du catalogue produits.

\subsection{Méthode 1 : L'Analyse par Marge de Transformation (\textit{Margin Spread Mode})}
\begin{itemize}
    \item \textbf{Dans quel cas l'utilise-t-on ?} : Cette méthode est le mode nominal standard. Elle est activée uniquement lorsque le prix du produit final (cible) ET les prix de ses matières premières (feedstocks) sont disponibles de façon continue dans notre base.
    \item \textbf{Principe Mathématique} : Le modèle calcule le Spread de marge brute $S_t$ en soustrayant le coût d'achat pondéré des précurseurs du prix de vente cible :
    \begin{equation}
    S_t = P^{\text{Produit Fini}}_t - \sum_{i=1}^{N} \left( \beta_i \times P^{\text{Matière } i}_t \right)
    \end{equation}
    \item \textbf{Intérêt Commercial} : Il permet aux acheteurs de négocier les approvisionnements en fonction de la rentabilité réelle des usines. Si le spread est historiquement bas, les fabricants produisent à perte et vont réduire leur offre, ce qui annonce une remontée imminente des prix.
\end{itemize}

\subsection{Méthode 2 : L'Analyse en Prix Absolu (\textit{Absolute Price Mode})}
\begin{itemize}
    \item \textbf{Dans quel cas l'utilise-t-on ?} : Utilisée pour les produits qui ont des données de prix fiables mais pour lesquels il n'y a pas d'intrants matières premières définis dans la chaîne de valeur (ex. le PTA ou le MEK), ou dont les matières premières ne sont pas disponibles sur OilChem.
    \item \textbf{Principe Mathématique} : L'écart-type, le RSI et les Bandes de Bollinger ne sont plus calculés sur une marge, mais directement sur la série de prix absolue du produit lui-même. Le spread théorique devient équivalent au prix :
    \begin{equation}
    S_t = P^{\text{Produit Fini}}_t
    \end{equation}
\end{itemize}

\subsection{Méthode 3 : L'Indexation sur Coût Matière Première (\textit{Raw Material Index Mode})}
\begin{itemize}
    \item \textbf{Dans quel cas l'utilise-t-on ?} : Utilisée pour les produits dont nous ne parvenons pas à obtenir les données de prix finis (produits de niche non cotés, ou bloqués par des restrictions d'accès OilChem), mais dont nous possédons le prix de la matière première principale.
    \item \textbf{Principe Mathématique} : Le modèle déplace le calcul des indicateurs techniques et la prévision Ridge sur le cours de la matière première principale ($P^{\text{Matière Principal}}_t$). Le spread de référence devient :
    \begin{equation}
    S_t = P^{\text{Matière Principal}}_t
    \end{equation}
    \item \textbf{Intérêt Commercial} : Cette méthode donne du **Pricing Power** (pouvoir de fixation des prix) à nos commerciaux (\textit{Sales}). Puisque la matière première représente 70\% à 85\% du coût d'un produit chimique, une hausse rapide de l'index matière première indique aux commerciaux qu'ils doivent immédiatement augmenter leurs prix de vente clients pour protéger la marge de DAXX EU, avant même que les prix du produit fini ne soient publiés.
\end{itemize}

\subsection{Cas Particulier : Le Mode Données Limitées (\textit{Low-Data Mode})}
Si un produit présente trop de valeurs manquantes (gaps prolongés) empêchant le calcul du RSI ou des Bandes de Bollinger (qui nécessitent un flux continu), le système bascule sur un mode simplifié. Il désactive les indicateurs techniques complexes et affiche un indicateur de \textbf{Momentum simple} (comparaison de prix à court et moyen terme sur 30j et 90j) combiné avec la saisonnalité historique.

---

\section{Modélisation et Notations du Profil de Risque (Financial Risk Profile)}
Cette section explique les indicateurs affichés sous le panneau **Financial Risk Profile** du volet latéral droit. Ces mesures évaluent les risques de volatilité financière et les seuils de pertes maximaux tolérés pour les transactions physiques de DAXX EU.

\subsection{1. La Volatilité Annualisée ($\sigma_{\text{annuelle}}$)}
La volatilité mesure l'instabilité ou la dispersion des rendements du prix ou du spread.
\begin{itemize}
    \item \textbf{Calcul} : Elle est calculée sur les 60 derniers jours de cotation. On calcule d'abord le rendement logarithmique quotidien $R_t = \ln(P_t / P_{t-1})$. Puis on calcule l'écart-type ($\text{std}$) de ces rendements et on l'annualise (en supposant 252 jours de trading par an) :
    \begin{equation}
    \sigma_{\text{annuelle}} = \text{std}(R_{t-59 \dots t}) \times \sqrt{252} \times 100
    \end{equation}
    \item \textbf{Interprétation} :
    \begin{itemize}
        \item \textbf{Risque Faible} : Volatilité $< 10\%$.
        \item \textbf{Risque Moyen} : Volatilité entre $10\%$ et $20\%$.
        \item \textbf{Risque Élevé} : Volatilité $> 20\%$. Un marché très volatil nécessite des marges de négociation plus larges.
    \end{itemize}
\end{itemize}

\subsection{2. La Value-at-Risk à 10 jours (VaR 95\%)}
La Value-at-Risk est une mesure standardisée en finance de marché pour estimer la perte maximale potentielle sur un horizon de temps donné avec un certain niveau de confiance.
\begin{itemize}
    \item \textbf{Calcul} : Le système extrait l'historique des variations absolues de prix sur un pas de temps de 10 jours : $\Delta P_{10d} = P_{\tau} - P_{\tau-10}$. Il calcule ensuite le 95ème centile de cette distribution de variations :
    \begin{equation}
    VaR_{10d, 95\%} = \text{Percentile}_{95\%}(\Delta P_{10d})
    \end{equation}
    \item \textbf{Interprétation} : Si la VaR est de 250 \yen, cela signifie que, dans \textbf{95\% des cas}, la hausse ou la variation maximale du prix sur les 10 prochains jours ne dépassera pas 250 \yen{} par tonne. C'est l'indicateur clé pour fixer les budgets de couverture de risques d'achats.
\end{itemize}

---

\section{Indicateurs de Performance du Modèle (Model Accuracy \& Value)}
Ce panneau valide la confiance que l'utilisateur peut accorder aux prédictions et chiffre l'économie réelle dégagée par l'outil.

\subsection{1. L'Erreur de Prévision (MAE \& MAPE)}
\begin{itemize}
    \item \textbf{MAE (Mean Absolute Error)} : Représente la distance absolue moyenne entre les prévisions à 14 jours et les prix réels observés.
    \item \textbf{MAPE (Mean Absolute Percentage Error)} : Représente le pourcentage d'erreur relatif de prévision.
    \begin{equation}
    \text{MAPE} = \frac{1}{M} \sum_{j=1}^{M} \left| \frac{\text{Prédiction}_j - \text{Prix Réel}_j}{\text{Prix Réel}_j} \right| \times 100
    \end{equation}
\end{itemize}

\subsection{2. La Précision Directionnelle (Model Accuracy / Signal Precision)}
\begin{itemize}
    \item \textbf{De quoi s'agit-il ?} : Pourcentage de fois où le modèle a correctement anticipé le sens d'évolution du marché (si le modèle a prédit une hausse et que le prix a effectivement monté, ou vice-versa).
    \item \textbf{Interprétation} : Une précision de 94\% indique que le modèle s'est orienté dans la bonne direction du marché dans 94 cas sur 100 historiques.
\end{itemize}

\subsection{3. L'Économie Financière (Financial Value / Savings)}
\begin{itemize}
    \item \textbf{De quoi s'agit-il ?} : Quantifie l'avantage d'une politique d'achat sélective dictée par nos signaux (acheter plus en état `FAVORABLE`, suspendre les achats en état `UNFAVORABLE`) par rapport à des achats passifs et continus.
    \item \textbf{Interprétation} : Si le tableau indique une économie de 250 \yen{} par tonne, cela signifie que sur l'historique testé, l'équipe d'achats a économisé en moyenne 250 \yen{} sur chaque tonne achetée grâce aux alertes de l'outil.
\end{itemize}

---

\newpage
\section{Démonstration Concrète : Cas du Butyl Acétate (华东)}
Voici une simulation mathématique complète pas à pas basée sur les équations de notre moteur de décision pour le **Butyl Acétate (Chine de l'Est)**.

\subsection{Données de Départ de la simulation}
Supposons les cours suivants enregistrés le jour $t$ dans la base de données :
\begin{itemize}
    \item Prix cible ($P^{\text{Cible}}_t$) du Butyl Acétate = \textbf{6 150 \yen / tonne}
    \item Prix du n-Butanol ($P^{\text{Butanol}}_t$) = \textbf{5 800 \yen / tonne}
    \item Prix de l'Acide Acétique ($P^{\text{Acide}}_t$) = \textbf{2 880 \yen / tonne}
\end{itemize}

\subsection{Étape 1 : Calcul du Coût Matières}
Coefficients techniques : $\beta_1 = 0.65$ (n-Butanol) et $\beta_2 = 0.35$ (Acide Acétique).
\begin{align}
\text{Coût Matières} &= (0.65 \times 5800) + (0.35 \times 2880) \\
\text{Coût Matières} &= 3770 + 1008 = \mathbf{4778\text{ \yen / tonne}}
\end{align}

\subsection{Étape 2 : Calcul du Spread}
\begin{align}
\text{Spread } S_t &= P^{\text{Cible}}_t - \text{Coût Matières} \\
S_t &= 6150 - 4778 = \mathbf{1372\text{ \yen / tonne}}
\end{align}

\subsection{Étape 3 : Position par rapport aux Bandes de Bollinger}
Pour les 20 derniers jours, les calculs statistiques de l'historique des spreads ont donné les résultats suivants :
\begin{itemize}
    \item Moyenne mobile simple $\mu_t = 1200\text{ \yen}$
    \item Écart-type local $\sigma_t = 80\text{ \yen}$
\end{itemize}
Calcul des bornes :
\begin{align}
\text{Bande Supérieure} &= 1200 + (2 \times 80) = \mathbf{1360\text{ \yen}} \\
\text{Bande Inférieure} &= 1200 - (2 \times 80) = \mathbf{1040\text{ \yen}}
\end{align}
Le spread actuel ($1372\text{ \yen}$) est supérieur à la Bande Supérieure ($1360\text{ \yen}$).

\subsection{Étape 4 : Calcul du RSI et Signal Final}
Le RSI lissé est calculé à **72.5** (surachat). 
Puisque le spread dépasse la Bande Supérieure et que le RSI est supérieur à 68, l'état de marché affecté est **`UNFAVORABLE`**.

\end{document}
